\section{Introduction}
\label{sec:introduction}

RNA molecules can be grouped into families based on shared characteristics such as structural patterns, sequence similarities, and functional roles. This allows researchers to study them in a more organized and meaningful way.  
To capture these similarities, a technique known as \mra{} (\mraabv{}) is often employed, in which several RNA sequences from the same family are aligned together.  
From this alignment, a model can be constructed that encodes the common features of the family. This model can then be used as a reference to evaluate new RNA reads, providing a measure of how likely it is that a given sequence belongs to the same RNA family.  

The chosen model for representing RNA families has often been a \smm{} (\smmabv{}), which is capable of describing many RNA molecules in terms of both their raw symbolic strings and certain structural configurations that the molecules can take.  
This model has the same expressive power as a \srg{} (\srgabv{}), meaning it can capture a wide range of sequence and structural patterns.  
However, it is limited in that it cannot represent more complex structures.  
These more intricate configurations are referred to as \sss{} (\sssabv{}), and they require more powerful modeling techniques to be fully captured.  
To address this issue, a \cm (\cmabv) is often used to model these structures. The \cmabv is a form of restricted \scfg (\scfgabv) which captures a subset of the
language of nested strings. These CM's are often represented as tree structures consisting of nodes. The nodes of the tree structure
either serve structural purposes, by adding in bifurcation to the structure, or contain states equipped with transition and emission probabilities.
% \br{verify this, it is rather apparent, but still find a source.}

% As we are primarily concerned with secondary structural properties, we choose to use the latter technique when comparing different \rnaf{}.  
% We hypothesize that \rnaf{} themselves can be examined for similarity, by comparing the similarity of the models used to represent them.  
% We propose a modified tree edit distance for examining the similarity of different \cm's which compares both structure and \lo transitions.
We hypothesize that the structural similarity of these CM trees can be used as a proxy of RNA Family similarity. To this 
end, we have developed an algorithm that utilizes the Zhang-Shasha Tree Edit distance algorithm, augmented with node to node comparison that uses an HMM Distance comparison 
function. This algorithm takes as input two CM's and measures their structural similarity.
The description of these algorithms is present in Section~\ref{sec:formalism} and our concrete implementation is described in Section~\ref{sec:implementation}.

To validate the efficacy of our approach we compare our distance calculation to the \clique{} of RNA families using PERMANOVA.
We present the evaluation of our approach in Section~\ref{sec:evaulation}.
Further, we discuss the weaknesses and potential benefits of our approach in Section~\ref{sec:discussion}.

% DOMINIC HAS VOLUNTEERED FOR INTRODUCTION EXPANSION
